%----------------------------------------------------------------------
% 結論與未來展望
%----------------------------------------------------------------------

\chapter{Conclusion and Future Work}\label{chap:conclusion_future_work}

\section{Conclusion}\label{sec:conclusion}

This thesis investigates consistency-based training for Multi-Token
Prediction (MTP) to improve the optimization of future-token prediction
heads. Two consistency objectives are introduced: Selective Distribution
Consistency (SDC), which selectively aligns the output distributions of
student heads with the first prediction head, and Latent Consistency
Matching (LCM), which aligns their latent representations. Both objectives
are incorporated into training without modifying the original MTP
architecture.

Experimental results on the 1.3B-parameter model demonstrate that SDC+LCM
improves future-token prediction and draft-token acceptance. Compared with
the baseline, Oracle PPL is reduced across all four prediction heads, with
Head 4 achieving the largest improvement of 5.89\%. The Average Draft Tokens
Accepted increases from 0.7010 to 0.7302, corresponding to an improvement
of 4.17\%. In addition, Generative PPL decreases from 39.5580 to 18.7180,
representing a reduction of 52.68\%. Zero-shot evaluation further shows
improvements across multiple unseen datasets.

The ablation study shows that the strengths of SDC and LCM, as well as the
confidence threshold, should be carefully balanced. Lower Generative PPL is
often accompanied by lower entropy and higher repetition, indicating a
trade-off with generation diversity. Overall, the proposed consistency-based
training framework improves the prediction quality of MTP heads, particularly
for more distant future-token predictions, while also increasing
draft-token acceptance without modifying the model architecture.

\section{Future Work}\label{sec:future_work}

Several directions can be explored in future work. First, the proposed
consistency objectives can be evaluated with larger models and larger
training budgets. Since the ablation study in this work was conducted
using models trained on only 1 billion tokens, larger-scale experiments
may provide a more comprehensive understanding of the effects of SDC,
LCM, and their associated hyperparameters.

Second, future studies can investigate the application of SDC and LCM
during supervised fine-tuning (SFT). This would help determine whether
the proposed consistency objectives remain effective when adapted to
downstream instruction-following and generation tasks.

Third, future work should evaluate the actual inference speed-up of the
proposed method using a complete self-speculative decoding system.
Although this work demonstrates an increase in Average Draft Tokens
Accepted, practical acceleration also depends on factors such as
verification overhead, memory access, and hardware utilization.
Therefore, wall-clock latency and end-to-end decoding throughput should
be measured to determine the practical efficiency gains of the proposed
approach.

Finally, the trade-off between consistency and generation diversity
deserves further investigation. Future work may explore alternative
regularization strategies or adaptive consistency weighting to preserve
the improvements in prediction quality while mitigating the increase
in repetitive generation.
